% ============ 第2章 §2.9 H^1与积分表示(原书 pp.147-149) ============
\section{\texorpdfstring{$H^1$ 与积分表示}{H1 与积分表示}}\label{sec:2.9}

2.4 节证明了属于 $h^{1}$ 的函数可以用 Poisson--Stieltjes 积分表示来刻画.本节证明属于 $\operatorname{Re}H^{1}$ 的函数可以用 Poisson 积分表示,即 $\mu$ 关于 Lebesgue 测度绝对连续.因此 $\operatorname{Re}H^{1} \subset h^{1}$.

\begin{theorem}\label{thm:2.9.1}
若 $f \in H^{1}$,则
\begin{equation*}
f(z) = \frac{1}{2\pi}\int_{-\pi}^{\pi}P(r, \theta - t)f(e^{it})\dd t.
\end{equation*}
\end{theorem}

\begin{proof}
当 $f \in H^{1}$ 时,若 $0 < \rho < 1$,则
\begin{equation*}
f(\rho z) = \frac{1}{2\pi}\int_{-\pi}^{\pi}P(r, \theta - t)f(\rho e^{it})\dd t,
\end{equation*}
且由定理 2.7.2,
\begin{equation*}
\int_{-\pi}^{\pi}|f(\rho e^{it}) - f(e^{it})|\dd t \rightarrow 0 \quad (\rho \rightarrow 1).
\end{equation*}
令
\begin{equation*}
F(z) = \frac{1}{2\pi}\int_{-\pi}^{\pi}P(r, \theta - t)f(e^{it})\dd t,
\end{equation*}
则 $f(\rho z) \to F(z)$,故 $F(z) = f(z)$.
\end{proof}

\begin{corollary}\label{cor:2.9.2}
若 $u \in \operatorname{Re}H^{1}$,则 $u$ 可以用 Poisson 积分表示.
\end{corollary}

对 $\phi \in L^{1}$,
\begin{equation*}
\hat{\phi}(n) = \frac{1}{2\pi}\int_{-\pi}^{\pi}\phi(t)e^{-int}\dd t \quad (n = 0, \pm 1, \pm 2, \ldots)
\end{equation*}
称为 $\phi$ 的 Fourier 系数.$\phi \in L^{1}$ 的 Poisson 积分在 $D$ 上调和,但何时正则?定理~\ref{thm:2.9.1} 表明当 $\phi$ 是 $H^{1}$ 的函数的边界值时正是如此.下面的命题~\ref{prop:2.9.3} 用 $\phi$ 的 Fourier 系数的语言回答这个问题.这将在下面的 2.10 节中用于 $H^{1}$ 的边界值的刻画.

\begin{proposition}\label{prop:2.9.3}
$\phi \in L^{1}$ 且 $\hat{\phi}(n) = 0$($n < 0$)时,令
\begin{equation*}
f(z) = \frac{1}{2\pi}\int_{-\pi}^{\pi}P(r, \theta - t)\phi(e^{it})\dd t \quad (z = re^{i\theta}),
\end{equation*}
则 $f \in H^{1}$ 且 $f(e^{i\theta}) = \phi(e^{i\theta})$ a.e. $\theta$.
\end{proposition}

\begin{proof}
由 $\hat{\phi}(n) = 0$($n < 0$)与 Poisson 核的展开
\begin{equation*}
P(r, \theta - t) = 1 + \sum_{n = 1}^{\infty}\left(z^{n}e^{-int} + \bar{z}^{n}e^{int}\right),
\end{equation*}
可以证明 $f(z) = \sum_{n=0}^{\infty}\hat{\phi}(n)z^{n}$($z \in D$).$f \in H(D)$ 由 $\phi \in L^{1}$ 是显然的,但由命题~\ref{prop:2.2.2} 的 (3),
\begin{equation*}
M_{1}(r, f) \leq \frac{1}{2\pi}\int_{-\pi}^{\pi}|\phi(e^{it})|\dd t,
\end{equation*}
所以 $f \in H^{1}$.由定理~\ref{thm:2.3.2},$f(e^{i\theta}) = \phi(e^{i\theta})$ a.e. $\theta$.
\end{proof}

若 $f \in H^{1}$,则
\begin{equation*}
f(z) = \frac{1}{2\pi}\int_{-\pi}^{\pi}C(r, \theta - t)f(e^{it})\dd t.
\end{equation*}
这可以与定理~\ref{thm:2.9.1} 同样地证明,但 2.12 节将给出另一个证明.
